\section{Conclusion}
This project compares five neural network configurations through a common
Fashion-MNIST pipeline. The data is split into fixed training, validation and
test partitions, normalized using training statistics, and presented as
vectors, image grids, row sequences or patch sequences. All models use the
same ten-epoch training budget and validation-loss checkpoint rule.

The CNN gives the strongest recorded result: 91.94\% test accuracy and
0.9189 macro-F1. It makes 806 errors on 10,000 images, 344 fewer than the MLP.
Its improvement appears across all ten class F1 scores, although Shirt remains
difficult, with 71.8\% recall. Most remaining CNN errors occur among
T-shirt/top, Pullover, Coat and Shirt.

MLP and LSTM obtain nearly equal accuracy at 88.50\% and 88.48\%. Their
two-image difference is too small to support a general ranking, and their
individual predictions differ substantially. Transformer reaches 85.55\%,
above the Linear baseline at 84.26\%, but below the other nonlinear models
under this configuration and training budget.

For this assignment, CNN is the preferred configuration when test accuracy
is the main criterion. MLP is a useful alternative when the recorded forward
cost matters: it trades some accuracy for a substantially shorter batch
forward time. LSTM is smaller than MLP in parameter count, but that smaller
count does not translate into faster measured inference.

The broader lesson is that a more elaborate architecture does not guarantee
a better result under a fixed experimental budget. Input structure, capacity,
optimization and runtime must be considered together. Repeated seeds and
controlled follow-up experiments are needed before these single-run findings
can be extended into stronger claims.

\vspace{8pt}
\begin{table}[H]\centering
\caption{Main findings and the evidence that supports them.}
\begin{tabularx}{\linewidth}{lX}\toprule
Finding & Evidence in this study\\\midrule
CNN leads in predictive quality &
91.94\% accuracy, 0.9189 macro-F1 and the highest F1 in each class.\\
MLP and LSTM are nearly tied &
88.50\% versus 88.48\% accuracy; no repeated-seed estimate.\\
Upper-body garments remain difficult &
\UpperErrors{} of the CNN's 806 errors occur within four upper-body labels.\\
Parameter count does not rank runtime &
LSTM and Transformer are smaller than MLP but have longer forward times.\\
Broader conclusions need more runs &
One seed, different capacities and separate timing sessions.\\\bottomrule
\end{tabularx}
\end{table}
